\dapso{Với 3,0 L E10 xe đi được khoảng \SI{144.8}{\kilo\metre}; cần đổ thêm khoảng \SI{0.73}{\liter} E10 để đi đủ \SI{180}{\kilo\metre}.}
\begin{loigiaiBT}
Thể tích ethanol:
\[
V_{\ce{C2H5OH}}=3.0\times10\%=\SI{0.30}{\liter}=\SI{300}{\milli\liter}.
\]
\[
m_{\ce{C2H5OH}}=300\times0.789=\SI{236.7}{\gram},\qquad
n_{\ce{C2H5OH}}=\frac{236.7}{46}\approx\SI{5.146}{\mole}.
\]
Phần xăng truyền thống có thể tích \SI{2.70}{\liter}, khối lượng:
\[
m=2700\times0.7=\SI{1890}{\gram}.
\]
Đặt $n_{\ce{C8H18}}=3x$, $n_{\ce{C9H20}}=4x$:
\[
3x\times114+4x\times128=1890\Rightarrow x\approx2.213.
\]
Suy ra:
\[
n_{\ce{C8H18}}\approx\SI{6.639}{\mole},\qquad
n_{\ce{C9H20}}\approx\SI{8.852}{\mole}.
\]
Tổng nhiệt lượng tỏa ra:
\[
Q\approx5.146\times1375+6.639\times5095+8.852\times6115
\approx\SI{95036}{\kilo\joule}.
\]
Công cơ học có ích:
\[
Q_{\text{ích}}=95036\times0.32\approx\SI{30411}{\kilo\joule}.
\]
Quãng đường đi được:
\[
s=\frac{30411}{210}\approx\SI{144.8}{\kilo\metre}.
\]
Do $144.8<180$, người đó phải đổ thêm xăng. Vì cùng loại E10, thể tích cần thêm tỉ lệ với quãng đường:
\[
V_{\text{thêm}}=3.0\left(\frac{180}{144.8}-1\right)\approx\SI{0.73}{\liter}.
\]
\textbf{Vậy,} người đó cần đổ thêm xăng; lượng E10 tối thiểu cần thêm khoảng \SI{0.73}{\liter}.
\end{loigiaiBT}
